% ============================================================
%  LEMBAR SOAL  --  FISIKA  --  SSU-ITB 2025  --  Paket 2
%  Set 1-6 (nomor 1--90) | PG Biasa
% ============================================================

\documentclass[11pt]{article}

\usepackage[utf8]{inputenc}
\usepackage[T1]{fontenc}
\usepackage[a4paper,margin=2.5cm,top=2.8cm]{geometry}
\usepackage{amsmath,amssymb}
\usepackage{graphicx}
\usepackage{enumitem}
\usepackage{multicol}
\usepackage{xcolor}
\usepackage[most]{tcolorbox}
\usepackage{fancyhdr}
\usepackage{booktabs}
\usepackage{icomma}
\usepackage{array}

\pagestyle{fancy}\fancyhf{}
\lhead{\small SSU-ITB 2025 --- Fisika}
\rhead{\small Paket 2}
\cfoot{\small\thepage}
\renewcommand{\headrulewidth}{0.4pt}

\newtcolorbox{soal}[1]{
  breakable,
  colback=white, colframe=black, colbacktitle=black,
  coltitle=white, fonttitle=\bfseries\sffamily,
  title=Nomor #1,
  boxrule=0.8pt, arc=0pt,
  left=3mm, right=3mm, top=2mm, bottom=2mm}

\newenvironment{pil}{%
  \begin{enumerate}[label=(\Alph*),leftmargin=*,itemsep=1pt,topsep=3pt]}%
  {\end{enumerate}}
\newenvironment{pilB}{%
  \par\smallskip\begin{multicols}{2}
  \begin{enumerate}[label=(\Alph*),leftmargin=*,itemsep=1pt,topsep=3pt]}%
  {\end{enumerate}\end{multicols}}

\newcommand{\gambar}[2][0.52\textwidth]{%
  \begin{center}\includegraphics[width=#1]{#2}\end{center}}

\linespread{1.05}

% ===================================================================
\begin{document}
\thispagestyle{empty}

\begin{center}
  \fbox{\parbox{0.6\textwidth}{\centering\bfseries\large LEMBAR SOAL}}
\end{center}
\vspace{0.5cm}

\begin{center}
  {\LARGE\bfseries FISIKA}
\end{center}
\vspace{0.5cm}

\begin{center}
  \begin{tabular}{@{\hspace{2cm}}ll@{\hspace{0.4cm}:@{\hspace{0.4cm}}}l}
    Jenjang      & & SMA / Sederajat               \\[0.3em]
    Hari/Tanggal & & \underline{\hspace{5cm}}      \\[0.3em]
    Waktu        & & \underline{\hspace{5cm}}      \\[0.3em]
    Paket Soal   & & 2                             \\
  \end{tabular}
\end{center}

\vspace{0.6cm}
\noindent\rule{\linewidth}{0.6pt}
\vspace{0.3cm}

\noindent\textbf{\underline{PETUNJUK UMUM}}
\begin{enumerate}[leftmargin=*,itemsep=4pt,topsep=4pt]
  \item Anda \textbf{tidak} diperbolehkan menggunakan sumber-sumber eksternal,
        termasuk internet, \textit{large language model} seperti ChatGPT, Claude,
        LLaMA, Gemini, Meta AI, dan sebagainya. Anda sangat dianjurkan untuk
        mencoba mengerjakan sendiri terlebih dahulu karena evaluasi ini dibuat
        untuk \textbf{mengukur} tingkat kepahaman; nilai $<70$ mengindikasikan
        \textbf{tidak lulus}, dan jika sudah melewati \textit{deadline} maka
        dianggap tidak mengerjakan yang berarti \textbf{tidak lulus}.

  \item Terdapat 2 tahap pengerjaan, yaitu tahap \textbf{Try Out} dan tahap
        \textbf{Review}. Try Out bersifat \textit{open book}; Anda diperbolehkan
        membuka buku apapun selama bersifat \textit{offline} dengan rentang waktu
        2 jam pada tanggal \underline{\hspace{3cm}}. Setelah itu memasuki
        tahap Review dengan \textit{schedule} rentang tanggal
        \underline{\hspace{3cm}}--\underline{\hspace{3cm}}, di mana Anda
        diharapkan mengerjakan ulang dengan disertakan penjelasan/pembelaan
        mengapa revisi jawaban adalah pilihan tersebut, dan tetap mencantumkan
        penjelasan pada soal yang walaupun pada penilaian sudah benar.

  \item Anda dianjurkan untuk mencantumkan caranya walaupun tidak termasuk
        penilaian, mengingat sifatnya adalah sebagai bahan pembelajaran.
        Mohon bertanggung jawab atas pekerjaan Anda sendiri.

  \item Bacalah setiap soal dengan teliti. Soal berjenis
        \textbf{Pilihan Ganda Biasa} --- 5 opsi (A--E), pilih
        1 jawaban paling tepat.

  \item Soal berjumlah \textbf{90 butir}: nomor 1--15 dari Set~1,
        nomor 16--30 dari Set~2, nomor 31--45 dari Set~3,
        nomor 46--60 dari Set~4, nomor 61--75 dari Set~5, dan
        nomor 76--90 dari Set~6 SSU-ITB 2025.
        Gunakan $g=10$ m/s$^2$ bila tidak disebutkan lain.

  \item Silakan bertanya apabila terdapat soal yang ambigu, rancu, atau
        kurang spesifik selama itu tidak menjurus kepada jawaban soal
        tersebut.

  \item Isilah detail informasi berikut.\\[0.3em]
        \textbf{Nama}\hspace{0.45cm};\enspace\underline{\hspace{8cm}}
\end{enumerate}

\vspace{0.6cm}
\noindent\textit{Dengan menandatangani kolom di bawah ini, saya menyatakan telah membaca,
memahami, dan menyetujui seluruh peraturan yang tercantum di atas.}

\vspace{0.6cm}
\noindent\fbox{\parbox{5cm}{\vspace{2.2cm}\centering\small Tanda Tangan Peserta\vspace{0.3cm}}}

\clearpage

% ===================================================================
%  SET 1  (nomor 1--15)
% ===================================================================

\begin{soal}{1}
Sebuah pesawat pengebom terbang dalam arah horizontal dengan laju konstan
$v_P=60$ m/s pada ketinggian $h=128$ m. Ketika tepat di atas titik $O$,
sebuah bom dilepaskan dan mengenai truk $T$ yang sedang bergerak di
jalan mendatar dengan kecepatan konstan. Pada saat bom dilepaskan, truk
berada pada jarak $x_0=100$ m dari $O$. Jika tinggi truk tersebut
adalah 3 m, maka laju truk $v_T$ adalah \ldots
\gambar[0.55\textwidth]{images/2/2-p2.png}
\begin{pilB}
  \item 65 m/s \item 60 m/s \item 40 m/s \item 50 m/s \item 30 m/s
\end{pilB}
\end{soal}

\begin{soal}{2}
Sebuah batu dilepaskan dari atas gedung dari kondisi diam. Pada saat yang
bersamaan dari ketinggian yang sama, sebuah bola dilemparkan secara
horizontal. Jika gesekan udara diabaikan, manakah pernyataan berikut ini
yang paling benar?
\begin{pil}
  \item Waktu sampai bola dan batu di tanah bergantung pada massanya masing-masing.
  \item Batu akan memiliki kelajuan yang lebih besar ketika sampai di tanah.
  \item Bola akan sampai di tanah terlebih dahulu.
  \item Batu akan sampai di tanah lebih dahulu.
  \item Batu dan bola sampai di tanah pada waktu bersamaan.
\end{pil}
\end{soal}

\begin{soal}{3}
Jika sebuah benda bermassa 5 kg memiliki energi kinetik sebesar 200 J
sesaat sebelum menumbuk tanah, maka dari ketinggian berapakah benda
tersebut dijatuhkan? Gunakan $g=10$ m/s$^2$.
\begin{pilB}
  \item 2 m \item 10 m \item 4 m \item 8 m \item 6 m
\end{pilB}
\end{soal}

\begin{soal}{4}
Dua buah balok dengan massa masing-masing $m=3{,}6$ kg dan $M=2$ kg
dihubungkan oleh sebuah tali yang melalui katrol seperti pada gambar.
Koefisien gesek statik dan kinetik antara dua permukaan balok adalah
$\mu_s=0{,}8$ dan $\mu_k=0{,}6$, tetapi lantai licin. Besar gaya
horizontal maksimum $F$ sehingga sistem masih berada dalam kesetimbangan
adalah \ldots\ N. Gunakan $g=10$ m/s$^2$.
\gambar[0.5\textwidth]{images/2/2-p4.png}
\begin{pilB}
  \item 61,60 \item 28,80 \item 86,40 \item 57,60 \item 64,80
\end{pilB}
\end{soal}

\begin{soal}{5}
Balok kecil $m_1=1$ kg dihubungkan dengan tali ke balok $m_2=4$ kg
melalui katrol licin tak bermassa. Sistem terletak pada lantai datar
seperti gambar. Koefisien gesek statis $m_1$ ke $m_2$ adalah 0,20 dan
koefisien gesek kinetik $m_2$ ke lantai adalah 0,30. Jika $m_2$ ditarik
dengan gaya mendatar $F=30$ N ke kiri, besar percepatan $m_1$ relatif
terhadap lantai selama bergerak di atas $m_2$ adalah \ldots
\gambar[0.5\textwidth]{images/2/2-p5.png}
\begin{pilB}
  \item 2,2 m/s$^2$ \item 2,4 m/s$^2$ \item 2,6 m/s$^2$
  \item 3,0 m/s$^2$ \item 4,4 m/s$^2$
\end{pilB}
\end{soal}

\begin{soal}{6}
Balok $m_1=1$ kg dihubungkan dengan tali ke balok $m_2=3$ kg melalui
katrol yang licin dan tak bermassa. Sistem pada lantai datar yang kasar.
Koefisien gesek kinetik antara $m_1$ ke $m_2$ dan antara $m_2$ ke lantai
berturut-turut adalah 0,25 dan 0,30. Jika $m_1$ ditarik dengan gaya
mendatar $F=35$ N, besar percepatan $m_2$ terhadap lantai selama balok
$m_1$ bergerak di atas $m_2$ adalah \ldots
\gambar[0.5\textwidth]{images/2/2-p6.png}
\begin{pilB}
  \item 5,375 m/s$^2$ \item 4,625 m/s$^2$ \item 3,850 m/s$^2$
  \item 3,500 m/s$^2$ \item 2,000 m/s$^2$
\end{pilB}
\end{soal}

\begin{soal}{7}
Persamaan gelombang mekanik di bawah ini yang berpanjang gelombang
0,75 m dan merambat dengan laju 15 m/s adalah:
\begin{pil}
  \item $y(t,x)=(0{,}40)\sin 2\pi\!\left(20t-\dfrac{4x}{3}+\dfrac{1}{3}\right)$ m
  \item $y(t,x)=(0{,}40)\sin 2\pi\!\left(20t-\dfrac{3x}{4}+\dfrac{1}{6}\right)$ m
  \item $y(t,x)=(0{,}40)\sin 2\pi\!\left(10t-\dfrac{4x}{3}+\dfrac{1}{5}\right)$ m
  \item $y(t,x)=(0{,}40)\sin 2\pi\!\left(10t-\dfrac{3x}{4}+\dfrac{1}{2}\right)$ m
  \item $y(t,x)=(0{,}40)\sin 2\pi\!\left(10t+\dfrac{4x}{3}+\dfrac{1}{2}\right)$ m
\end{pil}
\end{soal}

\begin{soal}{8}
Persamaan simpangan sebuah gelombang berjalan adalah
\[
  y(t,x) = 10\,\text{m}\,\sin\!\left(10\pi x + 15\pi t + \frac{\pi}{3}\right),
\]
dengan $t$ dalam sekon dan $x$ dalam meter. Pernyataan yang benar adalah \ldots
\begin{pil}
  \item Tidak ada jawaban yang benar.
  \item Gelombang merambat ke arah sumbu $x$ positif dengan panjang gelombang 0,2 m.
  \item Gelombang merambat ke arah sumbu $x$ positif dengan laju rambat 1,5 m/s.
  \item Gelombang merambat ke arah sumbu $x$ negatif dengan panjang gelombang 0,3 m.
  \item Gelombang merambat ke arah sumbu $x$ negatif dengan laju rambat 1,5 m/s.
\end{pil}
\end{soal}

\begin{soal}{9}
Gambar merupakan grafik yang menyatakan simpangan gelombang berjalan
pada $t=0$ yang merambat ke kanan. Jika frekuensi gelombang 5 Hz,
maka persamaan simpangan gelombang tersebut adalah (dalam SI) \ldots
\gambar[0.55\textwidth]{images/2/2-p8.png}
\begin{pil}
  \item $y(x,t)=2\cos(2\pi(5t-0{,}5x))$
  \item $y(x,t)=2\cos(2\pi(2t-x))$
  \item $y(x,t)=2\cos(2\pi(5t-x))$
  \item $y(x,t)=2\cos(2\pi(10t-2x))$
  \item $y(x,t)=2\cos(2\pi(10t-x))$
\end{pil}
\end{soal}

\begin{soal}{10}
Balok 2 kg pada gambar di bawah meluncur menuruni suatu lorong
melengkung yang licin dari keadaan diam pada ketinggian 3 m. Balok
kemudian bergerak pada permukaan horizontal $BC$ yang kasar dan berhenti
di $C$. Berapakah koefisien gesek antara balok dan permukaan horizontal?
\gambar[0.5\textwidth]{images/2/2-p9.png}
\begin{pilB}
  \item $\tfrac{1}{4}$ \item $\tfrac{1}{3}$ \item $\tfrac{1}{2}$
  \item $\tfrac{2}{5}$ \item $\tfrac{3}{5}$
\end{pilB}
\end{soal}

\begin{soal}{11}
Sebuah balok kecil bermassa 102 gram dilepaskan dari titik $A$. Balok
bergerak melewati lintasan lengkung $AB$ yang licin berbentuk seperempat
lingkaran. Setelah itu balok bergerak di lintasan datar $BC$ dengan
koefisien gesek kinetis 0,22 dan tepat berhenti di titik $C$. Rasio
antara jari-jari kelengkungan lintasan $AB$ dengan jarak lintasan datar
$BC$ adalah \ldots\ Gunakan $g=10$ m/s$^2$.
\gambar[0.5\textwidth]{images/2/2-p10.png}
\begin{pilB}
  \item 0,44 \item 0,11 \item 1,10 \item 0,22 \item 4,55
\end{pilB}
\end{soal}

\begin{soal}{12}
Sebuah balok meluncur pada bidang yang licin dari $P$ menuju $Q$.
Jika $g=10$ m/s$^2$, berapakah kelajuan balok di $P$ agar saat
mencapai titik $Q$ kelajuannya 7 m/s?
\gambar[0.5\textwidth]{images/2/2-p11.png}
\begin{pilB}
  \item 13 m/s \item 12 m/s \item 11 m/s \item 10 m/s \item 9 m/s
\end{pilB}
\end{soal}

\begin{soal}{13}
Sebuah balok yang diam bermassa 9,8 kg ditarik oleh gaya horizontal
konstan $F=10$ N. Jika permukaan lantai datar dan kasar dengan koefisien
gesek statik dan kinetik masing-masing 0,6 dan 0,4, maka percepatan
balok adalah \ldots\ m/s$^2$. Gunakan $g=10$ m/s$^2$.
\gambar[0.45\textwidth]{images/2/2-p12.png}
\begin{pilB}
  \item 3,12 \item 0,00 \item 0,88 \item 5,12 \item 4,00
\end{pilB}
\end{soal}

\begin{soal}{14}
Sebuah balok ditarik oleh gaya $F$ menaiki bidang miring. Massa balok
$m=5$ kg, sudut kemiringan $\theta=37^\circ$, koefisien gesek kinetik
$\mu_k=0{,}25$. Jika $F=45$ N sejajar bidang, besar percepatan balok
adalah \ldots
\gambar[0.5\textwidth]{images/2/2-p13.png}
\begin{pilB}
  \item 0,0 m/s$^2$ \item 1,0 m/s$^2$ \item 1,5 m/s$^2$
  \item 2,0 m/s$^2$ \item 3,0 m/s$^2$
\end{pilB}
\end{soal}

\begin{soal}{15}
Sebuah balok didorong oleh gaya $F$ \emph{horizontal} menaiki bidang
miring. Massa balok $m=5$ kg, sudut kemiringan $\theta=37^\circ$,
koefisien gesek kinetik $\mu_k=0{,}20$. Jika $F=75$ N, besar percepatan
balok adalah \ldots\ m/s$^2$.
\gambar[0.5\textwidth]{images/2/2-p14.png}
\begin{pilB}
  \item 3,6 \item 3,2 \item 3,0 \item 2,6 \item 2,4
\end{pilB}
\end{soal}

% ===================================================================
%  SET 2  (nomor 16--30)
% ===================================================================

\begin{soal}{16}
Sebuah partikel bermassa 500 gram mula-mula bergerak dengan kecepatan
6,5 m/s. Partikel tersebut kemudian ditarik oleh sebuah gaya konstan
$F$ sehingga dipercepat 2 m/s$^2$ selama 7 detik. Besar usaha yang
dilakukan oleh gaya $F$ selama dipercepat adalah \ldots\ Joule.
\begin{pilB}
  \item 5,12 \item 3,50 \item 105,6 \item 47,25 \item 94,50
\end{pilB}
\end{soal}

\begin{soal}{17}
Gaya $F_x$ bekerja pada balok $m=2$ kg yang berada pada bidang datar
licin seperti gambar. Kelajuan balok saat di $x=0$ adalah 4 m/s ke
sumbu $x$ positif. Berapakah kelajuan balok saat di $x=5$ m?
\gambar[0.5\textwidth]{images/2/2-p16.png}
\begin{pilB}
  \item 1,0 m/s \item 1,5 m/s \item 2,0 m/s \item 2,5 m/s \item 3,0 m/s
\end{pilB}
\end{soal}

\begin{soal}{18}
Sebuah balok bermassa $m=4$ kg ditarik dengan gaya $F=25$ N pada bidang
datar (arah $F$ membentuk sudut $37^\circ$ terhadap horizontal).
Koefisien gesek kinetik antara balok dan bidang adalah $\mu_k=0{,}20$.
Selama 4 detik, usaha yang diberikan oleh gaya $F$ ke balok adalah
\ldots\ Joule.
\begin{pilB}
  \item 150 \item 300 \item 450 \item 600 \item 750
\end{pilB}
\end{soal}

\begin{soal}{19}
Sebuah partikel bergerak dari pusat koordinat menuju arah sumbu $x$
positif dengan besar kecepatan 4 m/s selama 50 sekon. Kemudian tanpa
mengubah arah, besar kecepatannya menjadi 5 m/s untuk 20 sekon
berikutnya. Partikel tersebut kemudian berbalik arah dan selama 30 sekon
bergerak dengan besar kecepatan 6 m/s. Besar kecepatan rata-rata
partikel dalam 100 sekon tersebut adalah \ldots\ m/s.
\begin{pilB}
  \item 4,80 \item 0,03 \item 1,00 \item 1,20 \item 5,00
\end{pilB}
\end{soal}

\begin{soal}{20}
Sebuah bus berangkat menuju kota tertentu dengan kecepatan tetap
30 km/jam. Setelah tiba di kota tujuan, tanpa istirahat langsung kembali
ke kota asal dengan kecepatan tetap 60 km/jam. Kelajuan rata-rata bus
selama perjalanan tersebut adalah \ldots\ km/jam.
\begin{pilB}
  \item 35 \item 40 \item 45 \item 50 \item 55
\end{pilB}
\end{soal}

\begin{soal}{21}
Dari keadaan diam, sebuah bus dipercepat dengan besar percepatan
2 m/s$^2$ selama 10 sekon, lalu dengan kecepatan tetap selama 5 sekon,
kemudian diperlambat 4 m/s$^2$ sampai berhenti. Jika lintasan bus
adalah garis lurus, kelajuan rata-rata bus sejak dipercepat sampai
berhenti kembali adalah \ldots\ m/s.
\begin{pilB}
  \item 12,5 \item 15,0 \item 17,5 \item 20,0 \item 22,5
\end{pilB}
\end{soal}

\begin{soal}{22}
Jika
\[
  y(x,t) = (6\,\text{mm})\sin\!\left[2\pi\!\left(0{,}2x-\frac{5t}{3}\right)\right]
\]
mendeskripsikan gelombang transversal pada seutas tali, berapa lama
waktu minimum yang diperlukan sebuah titik pada tali untuk bergeser dari
$y=+3$ mm ke $y=-3$ mm?
\begin{pilB}
  \item 0,1 s \item 0,2 s \item 0,3 s \item 0,4 s \item 0,6 s
\end{pilB}
\end{soal}

\begin{soal}{23}
Persamaan simpangan sebuah gelombang berjalan transversal pada tali adalah
\[
  y(t,x)=(0{,}9\,\text{m})\sin\!\left(\frac{\pi x}{5}+8\pi t\right),
\]
dengan $t$ dalam sekon dan $x$ dalam meter. Laju transversal pada tali
di posisi $x=1{,}5$ m saat $t=0{,}5$ s adalah \ldots\ m/s.
\begin{pilB}
  \item 0,53 \item 13,29 \item 0,73 \item 18,29 \item 1,66
\end{pilB}
\end{soal}

\begin{soal}{24}
Gelombang merambat pada tali dengan persamaan $y=15\cos(4x-20t)$,
dengan $x$ dalam meter, $y$ dalam cm, dan $t$ dalam sekon. Kelajuan
getar tali pada posisi vertikal 12 cm adalah \ldots
\begin{pilB}
  \item 120 cm/s \item 150 cm/s \item 180 cm/s \item 200 cm/s \item 240 cm/s
\end{pilB}
\end{soal}

\begin{soal}{25}
Pada $t=0$ s, ambulans bergerak dengan kecepatan konstan 35 m/s menuju
suatu perempatan yang berjarak 1 km di depannya. Ambulans membunyikan
sirine dengan frekuensi 1000 Hz. Tepat saat 1 menit, frekuensi sirine
yang dideteksi oleh sebuah sensor yang dipasang di perempatan tersebut
adalah \ldots\ Hz. Cepat rambat bunyi di udara 340 m/s.
\begin{pilB}
  \item 906,67 \item 1000,00 \item 897,06 \item 1102,94 \item 1114,75
\end{pilB}
\end{soal}

\begin{soal}{26}
Sebuah sumber bunyi dengan frekuensi $f_s$ diletakkan di titik $x=0$.
Sumber bunyi tersebut kemudian berosilasi harmonik sepanjang sumbu $x$
dengan amplitudo osilasi $A_0$ dan frekuensi osilasi $f_0$. Sebuah
detektor bunyi diletakkan pada posisi $x=1{,}5A_0$. Manakah pernyataan
yang benar?
\begin{pil}
  \item frekuensi yang ditangkap oleh detektor saat sumber tepat berada di
        titik $x=A_0$ lebih besar dibanding saat sumber tepat di $x=-A_0$
  \item frekuensi yang ditangkap oleh detektor saat sumber tepat berada di
        titik $x=A_0$ dan $x=-A_0$ sama
  \item frekuensi yang ditangkap oleh detektor saat sumber tepat berada di
        titik $x=0$ dan sedang menuju $x=-A_0$ merupakan frekuensi maksimum
        yang terdeteksi
  \item frekuensi yang ditangkap oleh detektor saat sumber tepat berada di
        titik $x=0$ dan sedang menuju $x=A_0$ merupakan frekuensi minimum
        yang terdeteksi
  \item frekuensi yang ditangkap oleh detektor saat sumber tepat berada di
        titik $x=A_0$ lebih kecil dibanding saat sumber tepat di $x=-A_0$
\end{pil}
\end{soal}

\begin{soal}{27}
Sebuah sensor yang dipasang di halte mendeteksi frekuensi sirine mobil
ambulans yang melintas berubah dari 544 Hz menjadi 480 Hz. Laju bunyi
di udara saat itu 320 m/s. Jika ambulans melintas dengan kecepatan tetap,
besar kecepatannya adalah \ldots
\begin{pilB}
  \item 20 km/jam \item 30 km/jam \item 54 km/jam \item 72 km/jam \item 80 km/jam
\end{pilB}
\end{soal}

\begin{soal}{28}
Sebuah partikel dengan massa 7,6 gram bertumbukan dengan sebuah partikel
lain dengan massa 3,8 gram yang sedang dalam keadaan diam. Jika kedua
partikel menempel satu sama lain setelah tumbukan, maka perbandingan
antara energi yang hilang akibat tumbukan dengan energi kinetik awal
adalah \ldots
\begin{pilB}
  \item 0,50 \item 0,25 \item 0,33 \item 0,00 \item 0,67
\end{pilB}
\end{soal}

\begin{soal}{29}
Pada tumbukan dua buah benda, diketahui energi kinetik sistem kekal.
Manakah pernyataan yang benar mengenai momentum sistem setelah tumbukan?
\begin{pil}
  \item momentum sistem mungkin kekal
  \item momentum sistem berkurang setengahnya
  \item momentum sistem pasti kekal
  \item momentum sistem naik menjadi dua kali momentum sebelum tumbukan
  \item momentum sistem sama dengan energi kinetik sistem
\end{pil}
\end{soal}

\begin{soal}{30}
Sebuah benda 5 kg dengan kelajuan 4,0 m/s menabrak secara sentral benda
10 kg yang bergerak ke arahnya dengan kelajuan 3,0 m/s. Jika setelah
tumbukan benda 10 kg berhenti, maka koefisien restitusi tumbukan dua
benda itu adalah \ldots
\begin{pilB}
  \item $\tfrac{2}{3}$ \item $\tfrac{4}{7}$ \item $\tfrac{3}{7}$
  \item $\tfrac{1}{3}$ \item $\tfrac{2}{7}$
\end{pilB}
\end{soal}

% ===================================================================
%  SET 3  (nomor 31--45)
% ===================================================================

\begin{soal}{31}
Bola ditendang dari tanah dengan sudut elevasi $\theta$ terhadap tanah.
Bola jatuh di tanah pada jarak mendatar yang sama dengan empat kali
tinggi maksimum yang dicapai bola. Jika gesekan dengan udara diabaikan,
maka \ldots
\begin{pilB}
  \item $\theta=30^\circ$ \item $\theta=37^\circ$ \item $\theta=45^\circ$
  \item $\theta=53^\circ$ \item $\theta=60^\circ$
\end{pilB}
\end{soal}

\begin{soal}{32}
Sebuah proyektil ditembakkan dari tanah di $A$ ke arah benda $B$. Pada
saat yang sama, $B$ dilepaskan. Jarak $AC=45$ m dan tinggi $B$ adalah
60 m vertikal di atas $C$. Berapa laju awal proyektil saat ditembakkan
agar tepat mengenai $B$ saat berada 15 m di atas tanah?
\gambar[0.42\textwidth]{images/2/2-p27.png}
\begin{pilB}
  \item 20 m/s \item 25 m/s \item 30 m/s \item 40 m/s \item 45 m/s
\end{pilB}
\end{soal}

\begin{soal}{33}
Seorang pelempar akan melemparkan sebuah bola dengan laju 30 m/s dan
sudut elevasi $\alpha$ dengan $\tan\alpha=3/4$. Seorang penangkap
berada pada jarak 60 m dari pelempar. Jika $g=10$ m/s$^2$ dan penangkap
dapat berlari dengan kecepatan konstan, maka agar penangkap dapat
menangkap bola, dia harus berlari dengan laju minimum sebesar \ldots\ m/s.
\begin{pilB}
  \item 10,0 \item 5,66 \item 4,33 \item 7,33 \item 24,0
\end{pilB}
\end{soal}

\begin{soal}{34}
Balok bermassa 5 kg ditekan pada dinding dengan koefisien gesek 0,3
dengan gaya tertentu seperti gambar. Gaya minimum yang diperlukan untuk
menghentikan balok agar tidak jatuh adalah \ldots\ N.
Gunakan $g=10$ m/s$^2$.
\begin{pilB}
  \item 200,0 \item 142,9 \item 133,3 \item 166,7 \item 100,0
\end{pilB}
\end{soal}

\begin{soal}{35}
Untuk menjaga agar kotak kayu bermassa 10 kg tetap menempel diam pada
dinding tegak, kotak didorong dengan gaya $F$ yang membentuk sudut
$37^\circ$ terhadap garis datar. Koefisien gesek statis antara kotak
dan dinding 0,25. Batas nilai $F$ adalah \ldots
\begin{pil}
  \item $105\,\text{N}\leq F\leq 270\,\text{N}$
  \item $110\,\text{N}\leq F\leq 265\,\text{N}$
  \item $125\,\text{N}\leq F\leq 250\,\text{N}$
  \item $120\,\text{N}\leq F\leq 255\,\text{N}$
  \item $115\,\text{N}\leq F\leq 260\,\text{N}$
\end{pil}
\end{soal}

\begin{soal}{36}
Balok $B$ bermassa $M=50$ kg di atas lantai datar yang licin. Balok $A$
bermassa $m=10$ kg didorong dengan gaya mendatar $F$ sehingga kontak ke
balok $B$. Koefisien gesek statis antara $A$ dan $B$ adalah 0,40.
Berapa nilai minimum $F$ agar balok $A$ tidak mengalami slip terhadap
balok $B$?
\begin{pilB}
  \item 300 N \item 325 N \item 350 N \item 400 N \item 600 N
\end{pilB}
\end{soal}

\begin{soal}{37}
Pada gambar~(a), batang homogen yang beratnya $W_1$ dan panjangnya $L$
setimbang horizontal karena ujung kiri diengselkan ke dinding sedangkan
ujung kanan ditahan oleh kabel tipis takbermassa. Pada batang diletakkan
beban yang beratnya $W_2$ pada jarak $x$ dari dinding, dengan
$W_1+W_2=700$ N. Gambar~(b) menunjukkan besar tegangan $T$ pada kabel
terhadap rasio $x/L$. Maka $W_2=\ldots$
\begin{pilB}
  \item 100 N \item 200 N \item 300 N \item 400 N \item 500 N
\end{pilB}
\end{soal}

\begin{soal}{38}
Batang $AB$ yang homogen beratnya 200 N bertumpu ke lantai datar yang
kasar di $A$ sedangkan ujung $B$ terikat ke tali ringan. Di ujung $B$
digantung beban yang beratnya $W=600$ N. Saat itu batang setimbang
hampir tergelincir di $A$. Koefisien gesek statis antara $A$ dengan
lantai adalah \ldots\ ($\sin 37^\circ=0{,}6$ dan $\sin 53^\circ=0{,}8$).
\begin{pilB}
  \item 0,40 \item 0,50 \item 0,60 \item 0,75 \item 0,80
\end{pilB}
\end{soal}

\begin{soal}{39}
Sebuah papan iklan homogen dan persegi panjang sisinya $L=2{,}0$ m dan
beratnya 600 N. Papan ditahan oleh batang tak bermassa panjangnya 3,0 m
serta posisinya horizontal. Ujung kiri batang diengselkan ke dinding
sedangkan ujung kanan ditahan oleh kabel takbermassa seperti gambar.
Berapakah gaya tegangan pada kabel?
\gambar[0.4\textwidth]{images/2/2-p31.png}
\begin{pilB}
  \item 750 N \item 1000 N \item 1250 N \item 1500 N \item 2000 N
\end{pilB}
\end{soal}

\begin{soal}{40}
Suatu sumber bunyi memiliki daya akustik $1\,\mu$W. Jika intensitas
ambang pendengaran $10^{-12}$ W/m$^2$, maka tingkat intensitas pada
jarak 10 m dari sumber bunyi adalah \ldots\ dB.
\begin{pilB}
  \item 40 \item 45 \item 50 \item 60 \item 80
\end{pilB}
\end{soal}

\begin{soal}{41}
Seorang siswa mengukur taraf intensitas bunyi dari sebuah sumber. Pada
jarak 10 m dari sumber siswa mengukur 60 dB. Berapa taraf intensitas
bunyi pada jarak 100 m dari sumber bunyi tersebut?
\begin{pilB}
  \item 6 dB \item 10 dB \item 40 dB \item 45 dB \item 50 dB
\end{pilB}
\end{soal}

\begin{soal}{42}
Agar taraf intensitas bunyi dari sebuah sumber berkurang dari 57 dB
menjadi 47 dB, maka daya akustik sumber bunyi itu harus dikurangi
sebanyak \ldots
\begin{pilB}
  \item 10\% \item 20\% \item 50\% \item 80\% \item 90\%
\end{pilB}
\end{soal}

\begin{soal}{43}
Gelombang stasioner yang dihasilkan pada seutas tali dinyatakan oleh
$y=5\cos\!\left(\dfrac{\pi x}{3}\right)\sin(40\pi t)$, dengan $x$ dan
$y$ dalam cm dan $t$ dalam detik. Jarak antara dua simpul berurutan
adalah \ldots\ cm.
\begin{pilB}
  \item 3 \item 5 \item 20 \item $\pi$ \item 40
\end{pilB}
\end{soal}

\begin{soal}{44}
Seutas kawat yang tipis panjangnya 9 m dan bermassa 0,15 kg direntangkan
dengan kedua ujungnya terikat kuat. Gaya tegangan pada kawat itu 135 N.
Frekuensi minimum gelombang pada kawat adalah \ldots
\begin{pilB}
  \item 2 Hz \item 5 Hz \item 6 Hz \item 8 Hz \item 10 Hz
\end{pilB}
\end{soal}

\begin{soal}{45}
Gempa bumi menghasilkan dua jenis gelombang yang merambat di tanah, yaitu
gelombang longitudinal P dengan laju 8,5 km/det dan gelombang transversal
S dengan laju 5,5 km/det. Apabila kedua gelombang terdeteksi seismograf
berselisih 2 menit, maka jarak antara seismograf terhadap pusat gempa
itu kira-kira \ldots\ km.
\begin{pilB}
  \item 1870 \item 1680 \item 1020 \item 660 \item 360
\end{pilB}
\end{soal}

% ===================================================================
%  SET 4  (nomor 46--60)
% ===================================================================

\begin{soal}{46}
Dua buah partikel identik A dan B masing-masing sedang bergerak
melingkar. Nilai percepatan sentripetal terhadap jari-jari lintasan
dari kedua partikel dapat dilihat pada gambar di bawah. Rasio besar
kecepatan sudut A terhadap besar kecepatan sudut B adalah \ldots
\gambar[0.45\textwidth]{images/2/2-p34.png}
\begin{pilB}
  \item $2:1$ \item $\sqrt{2}:1$ \item $3:1$ \item $1:\sqrt{2}$ \item $1:2$
\end{pilB}
\end{soal}

\begin{soal}{47}
Sebuah benda yang massanya 8 kg bergerak secara beraturan dalam lintasan
melingkar dengan laju 5 m/s. Bila jari-jari lingkaran itu 1 m, maka
\ldots
\begin{enumerate}[label=\alph*.,leftmargin=*,itemsep=2pt,topsep=3pt]
  \item gaya sentripetalnya adalah 200 N
  \item waktu putarnya adalah 0,4 sekon
  \item vektor kecepatannya tidak tetap
  \item besar percepatan sentripetalnya 25 m/s$^2$
\end{enumerate}
Pernyataan yang benar adalah \ldots
\begin{pil}
  \item hanya a, b, dan c \item hanya a dan c \item hanya b dan d
  \item hanya d \item semua pernyataan benar
\end{pil}
\end{soal}

\begin{soal}{48}
Sebuah bola diikat dengan tali lalu diputar pada bidang vertikal. Massa
bola 300 g, radius lintasannya 75 cm. Jika bola berputar dengan laju
tetap 4,0 m/s, besarnya gaya tegangan pada tali saat bola di puncak
lintasannya adalah \ldots\ N.
\begin{pilB}
  \item 3,4 \item 3,6 \item 4,5 \item 6,4 \item 9,4
\end{pilB}
\end{soal}

\begin{soal}{49}
Seorang pemain menendang bola bermassa 0,40 kg. Bola meninggalkan kaki
pemain dengan laju awal 25 m/s. Jika antara bola dan kaki terjadi kontak
selama 8 ms, besar gaya rata-rata yang diberikan kaki ke bola selama
terjadinya kontak adalah \ldots\ N.
\begin{pilB}
  \item 950 \item 1150 \item 1250 \item 1500 \item 1750
\end{pilB}
\end{soal}

\begin{soal}{50}
Mobil mainan bermassa 5 kg mula-mula diam di titik asal. Kemudian pada
mobil mainan tersebut diberikan gaya mendatar $F_x$ yang berubah terhadap
waktu sesuai gambar di bawah. Kelajuan mobil mainan itu saat $t=9{,}0$ s
adalah \ldots
\gambar[0.5\textwidth]{images/2/2-p44.png}
\begin{pilB}
  \item 12 m/s \item $2\sqrt{5}$ m/s \item $5\sqrt{2}$ m/s
  \item $2\sqrt{3}$ m/s \item $3\sqrt{2}$ m/s
\end{pilB}
\end{soal}

\begin{soal}{51}
Balok bermassa 2,0 kg bergerak dengan kecepatan $(2{,}0\hat{\imath}-3{,}0\hat{\jmath})$
m/s dan balok bermassa 3,0 kg bergerak dengan kecepatan
$(1{,}0\hat{\imath}+6{,}0\hat{\jmath})$ m/s. Kecepatan pusat massa sistem
dua balok tersebut adalah \ldots\ m/s.
\begin{pil}
  \item $(7{,}0\hat{\imath}+12{,}0\hat{\jmath})$
  \item $(1{,}4\hat{\imath}+2{,}4\hat{\jmath})$
  \item $(7{,}0\hat{\imath}+2{,}4\hat{\jmath})$
  \item $(1{,}4\hat{\imath}-2{,}4\hat{\jmath})$
  \item $(-1{,}4\hat{\imath}+2{,}4\hat{\jmath})$
\end{pil}
\end{soal}

\begin{soal}{52}
Seorang anak laki-laki berdiri di timbangan yang berada di dalam lift.
Ketika lift bergerak naik jarum timbangan menunjuk angka 621 N, sedangkan
saat lift akan berhenti jarum timbangan menunjuk angka 459 N. Diketahui
$g=10$ m/s$^2$. Besar percepatan saat lift naik sama dengan besar
perlambatan saat lift akan berhenti. Massa anak laki-laki itu adalah
\ldots\ kg.
\begin{pilB}
  \item 50 kg \item 52 kg \item 54 kg \item 58 kg \item 60 kg
\end{pilB}
\end{soal}

\begin{soal}{53}
Seorang anak bermassa 40 kg, ketika ditimbang di dalam lift yang bergerak
vertikal massanya seolah-olah menjadi 60 kg. Hal itu berarti lift sedang
bergerak \ldots
\begin{pil}
  \item naik, dengan perlambatan 0,5 m/s$^2$
  \item turun, dengan perlambatan 0,5 m/s$^2$
  \item naik, dengan percepatan 1,0 m/s$^2$
  \item turun, dengan percepatan 1,0 m/s$^2$
  \item naik, dengan kecepatan tetap
\end{pil}
\end{soal}

\begin{soal}{54}
Perhatikan sistem pada gambar. Massa benda A adalah 10 kg dan massa
benda B adalah 5 kg. Jika benda A ditarik ke atas sehingga sistem
bergerak ke atas dengan percepatan konstan sebesar 2 m/s$^2$ dan
$g=10$ m/s$^2$, maka besar tegangan tali $T_1$ adalah \ldots
\begin{pilB}
  \item 30 N \item 90 N \item 10 N \item 180 N \item 60 N
\end{pilB}
\end{soal}

\begin{soal}{55}
Dua buah balok bermassa $m_1=5$ kg dan $m_2=8$ kg terikat pada tali
yang melalui sebuah katrol seperti gambar. Massa katrol, massa tali, dan
gesekan katrol diabaikan. Balok $m_1$ berada di atas meja horizontal
kasar dengan koefisien gesek statik dan kinetik 0,5 dan 0,4. Jika
$g=10$ m/s$^2$, maka besar gaya tegangan pada tali adalah \ldots\ N.
\gambar[0.42\textwidth]{images/2/2-p40.png}
\begin{pilB}
  \item 186,7 \item 14,0 \item 8,6 \item 18,5 \item 43,1
\end{pilB}
\end{soal}

\begin{soal}{56}
Tiga balok bermassa $m_A=30$ kg, $m_B=40$ kg, dan $m_C=10$ kg
dihubungkan dengan tali melalui katrol takbermassa seperti gambar.
Koefisien gesek kinetis antara balok A dengan meja adalah 0,20.
Besar gaya tegangan pada tali yang menghubungkan B ke C adalah \ldots\ N.
\gambar[0.4\textwidth]{images/2/2-p45.png}
\begin{pilB}
  \item 35 \item 40 \item 45 \item 225 \item 375
\end{pilB}
\end{soal}

\begin{soal}{57}
Kotak bermassa $m_2=3{,}5$ kg diam pada permukaan meja licin dan
dihubungkan ke kotak $m_1=1{,}5$ kg dan kotak $m_3=2{,}5$ kg dengan
tali ringan yang melalui katrol licin takbermassa. Jika $g=10$ m/s$^2$,
besar percepatan kotak $m_1$ adalah \ldots\ m/s$^2$.
\gambar[0.42\textwidth]{images/2/2-p46.png}
\begin{pilB}
  \item $5/4$ \item $4/3$ \item $3/2$ \item $2$ \item $1$
\end{pilB}
\end{soal}

\begin{soal}{58}
Proyektil ditembakkan dari tanah tegak lurus ke atas dengan laju awal
60 m/s. Gesekan udara diabaikan dan $g=10$ m/s$^2$. Kelajuan proyektil
saat mencapai posisi yang tingginya setengah dari tinggi maksimum adalah
\ldots\ m/s.
\begin{pilB}
  \item 30 \item 35 \item $30\sqrt{2}$ \item 40 \item 44
\end{pilB}
\end{soal}

\begin{soal}{59}
Jika bola jatuh bebas dari ketinggian tertentu, maka bola tersebut
menumbuk tanah dengan laju 30 m/s. Gesekan udara diabaikan. Jika bola
dari tempat yang sama dilemparkan vertikal ke bawah dengan laju awal
40 m/s, maka bola akan menumbuk tanah dengan kelajuan \ldots\ m/s.
\begin{pilB}
  \item 45 \item 50 \item 55 \item 60 \item 70
\end{pilB}
\end{soal}

\begin{soal}{60}
Dari ketinggian 120 m di atas tanah sebuah bola dilepaskan jatuh bebas.
Pada saat yang sama, bola lain dilemparkan vertikal ke atas dari tanah
dengan laju awal 60 m/s. Kedua bola itu berpapasan pada ketinggian
\ldots\ m di atas tanah.
\begin{pilB}
  \item 100 \item 90 \item 80 \item 70 \item 80
\end{pilB}
\end{soal}

% ===================================================================
%  SET 5  (nomor 61--75)
% ===================================================================

\begin{soal}{61}
Grafik di bawah mendeskripsikan kecepatan sebuah partikel yang bergerak
pada sumbu $x$. Gerak partikel dimulai dari titik asal ($x=0$ m) ke
arah sumbu $x$ positif. Kelajuan rata-rata partikel pada selang waktu
$t=0$ s sampai $t=9$ s adalah \ldots
\gambar[0.5\textwidth]{images/2/2-s5n1.png}
\begin{pilB}
  \item $28/9$ m/s \item $30/9$ m/s \item $36/9$ m/s
  \item $40/9$ m/s \item $44/9$ m/s
\end{pilB}
\end{soal}

\begin{soal}{62}
Partikel yang mula-mula rehat di titik asal lalu mengalami percepatan
searah sumbu $x$. Grafik antara besar percepatan terhadap waktu
ditunjukkan pada gambar. Pada selang waktu $t=0$ sampai $t=40$ s,
kelajuan rata-rata dan kecepatan rata-rata partikel secara berurutan
adalah \ldots
\gambar[0.5\textwidth]{images/2/2-s5n2.png}
\begin{pil}
  \item $+7{,}5000$ m/s, $0$ m/s
  \item $+7{,}8125$ m/s, $-0{,}3125$ m/s
  \item $+7{,}8125$ m/s, $+0{,}3125$ m/s
  \item $-0{,}3125$ m/s, $+7{,}8125$ m/s
  \item $+7{,}5000$ m/s, $-7{,}5000$ m/s
\end{pil}
\end{soal}

\begin{soal}{63}
Saat lampu lalu lintas menyala hijau, mobil sedan yang berhenti mulai
bergerak dengan percepatan tetap $3{,}2$ m/s$^2$. Pada saat bersamaan,
sebuah truk melintas di samping sedan dengan kecepatan tetap 20 m/s.
Perhatikan pernyataan berikut: (1)~sedan tepat menyusul truk pada
$t=12{,}5$ detik; (2)~sedan tepat menyusul truk saat bergerak sejauh
250 m; (3)~sedan tepat menyusul truk saat kecepatan sedan 20 m/s.
Pernyataan yang benar adalah \ldots
\begin{pilB}
  \item hanya (1) \item hanya (2) \item hanya (3)
  \item hanya (1) dan (2) \item (1), (2), dan (3) benar
\end{pilB}
\end{soal}

\begin{soal}{64}
Sebuah batu dilepaskan di tepi puncak sebuah gedung. Dua detik kemudian
di tempat yang sama sebuah bola dilemparkan dengan laju awal 30 m/s
tegak lurus ke bawah. Diketahui $g=10$ m/s$^2$. Jika batu dan bola
mencapai tanah pada saat yang bersamaan, tinggi gedung tersebut adalah
\ldots\ m.
\begin{pilB}
  \item 60 \item 70 \item 80 \item 90 \item 100
\end{pilB}
\end{soal}

\begin{soal}{65}
Tiga buah vektor gaya masing-masing $\vec{F}_1=(-2{,}0\hat{\imath}+2{,}0\hat{\jmath})$
N, $\vec{F}_2=(5{,}0\hat{\imath}-3{,}0\hat{\jmath})$ N, dan
$\vec{F}_3=(-45{,}0\hat{\imath})$ N bekerja pada sebuah benda. Jika
akibat ketiga gaya itu benda mengalami percepatan yang besarnya
$3{,}75$ m/s$^2$, maka massa benda itu adalah \ldots\ kg.
\begin{pilB}
  \item 10,5 \item 11,2 \item 12,3 \item 12,6 \item 13,2
\end{pilB}
\end{soal}

\begin{soal}{66}
Sebuah beban bermassa $m=50$ g digantung di dalam truk yang sedang rehat
di jalan datar. Saat truk dipercepat dengan percepatan $a$ ke kanan,
beban menyimpang ke kiri sampai tali mengapit sudut $\theta=37^\circ$.
Jika $g=10$ m/s$^2$, maka $a=\ldots$ m/s$^2$.
\gambar[0.45\textwidth]{images/2/2-p53.png}
\begin{pilB}
  \item 6,00 \item 6,25 \item 7,00 \item 7,50 \item 8,00
\end{pilB}
\end{soal}

\begin{soal}{67}
Tiga balok bermassa $m_1=2$ kg, $m_2=3$ kg, dan $m_3=4$ kg disusun pada
bidang datar licin, selanjutnya didorong oleh gaya $F=18$ N seperti
gambar. Besar gaya kontak $m_1$ ke $m_2$ serta gaya kontak $m_2$ ke
$m_3$ secara berurutan adalah \ldots
\gambar[0.45\textwidth]{images/2/2-p56.png}
\begin{pil}
  \item 4 N\,;\,8 N \item 4 N\,;\,10 N \item 14 N\,;\,10 N
  \item 16 N\,;\,8 N \item 14 N\,;\,8 N
\end{pil}
\end{soal}

\begin{soal}{68}
Empat balok pada lantai datar dihubungkan oleh tali ringan seperti gambar.
Koefisien gesek kinetis tiap balok ke lantai masing-masing 0,1. Jika
susunan itu ditarik dengan gaya mendatar $F=50$ N, maka gaya tegangan
pada tali nomor~(2) adalah \ldots\ N. Gunakan $g=10$ m/s$^2$. (Susunan
dari belakang ke depan: 10 kg --- 3 kg --- 5 kg --- 2 kg, lalu $F$.)
\begin{pilB}
  \item 32,5 \item 22,5 \item 19,5 \item 17,5 \item 15,5
\end{pilB}
\end{soal}

\begin{soal}{69}
Sebuah inti atom yang tidak stabil bermassa $17{,}0\times10^{-27}$ kg
mula-mula rehat. Sesaat kemudian inti tersebut meluruh menjadi tiga
partikel. Partikel pertama bermassa $5{,}0\times10^{-27}$ kg bergerak
pada arah sumbu $y$ positif dengan kelajuan $6{,}0\times10^{6}$ m/s.
Partikel kedua bermassa $8{,}4\times10^{-27}$ kg bergerak pada arah
sumbu $x$ negatif dengan kelajuan $4{,}0\times10^{6}$ m/s. Kecepatan
partikel ketiga adalah \ldots
\begin{pil}
  \item $(9{,}33\hat{\imath}+8{,}33\hat{\jmath})\times10^6$ m/s
  \item $(-9{,}33\hat{\imath}+8{,}33\hat{\jmath})\times10^6$ m/s
  \item $(-9{,}33\hat{\imath}-8{,}33\hat{\jmath})\times10^6$ m/s
  \item $(9{,}33\hat{\imath}-8{,}33\hat{\jmath})\times10^6$ m/s
  \item $(9{,}33\hat{\imath}-8{,}33\hat{k})\times10^6$ m/s
\end{pil}
\end{soal}

\begin{soal}{70}
Sebuah balok bermassa 3,0 kg bergerak dengan kecepatan awal
$5{,}0\hat{\imath}$ m/s, bertumbukan dengan balok bermassa 2,0 kg yang
kecepatan awalnya $-3{,}0\hat{\jmath}$ m/s. Setelah bertumbukan, kedua
balok menjadi satu dan bergerak dengan kelajuan \ldots\ m/s.
\begin{pilB}
  \item 3,10 \item 3,23 \item 3,56 \item 3,65 \item 3,75
\end{pilB}
\end{soal}

\begin{soal}{71}
Gelombang sinus yang merambat di kawat yang ditegangkan horizontal
dinyatakan melalui persamaan $y=0{,}20\sin(0{,}75\pi x+18\pi t)$,
dengan $x$ dan $y$ dalam meter dan $t$ dalam detik. Jika kerapatan
linier kawat 0,25 kg/m, besar gaya tegangan pada kawat adalah \ldots\ N.
\begin{pilB}
  \item 72 \item 144 \item 164 \item 212 \item 360
\end{pilB}
\end{soal}

\begin{soal}{72}
Dua tali yang kerapatannya $\mu_1=0{,}10$ kg/m dan $\mu_2=0{,}20$ kg/m
disambungkan pada suatu titik lalu ditegangkan. Rasio panjang gelombang
pada kedua kawat ($\lambda_1:\lambda_2$) adalah \ldots
\begin{pilB}
  \item $1:2$ \item $1:4$ \item $1:\sqrt{2}$ \item $2:1$ \item $1:2\sqrt{2}$
\end{pilB}
\end{soal}

\begin{soal}{73}
Ujung seutas tali dihubungkan ke osilator di $P$, ujung lainnya ditahan
di $Q$, lalu tali ditegangkan oleh beban bermassa $m$. Jarak $PQ$
adalah $L=1{,}20$ m, osilator berfrekuensi 150 Hz, sedangkan kerapatan
jenis tali $\mu=2{,}0$ g/m. Jika gelombang berdiri yang timbul di tali
memuat empat perut, maka $m=\ldots$ kg.
\begin{pilB}
  \item 1,62 \item 1,75 \item 1,84 \item 2,12 \item 2,25
\end{pilB}
\end{soal}

\begin{soal}{74}
Jika intensitas gelombang gempa yang terukur pada jarak 100 km dari
sumbernya adalah $1\times10^{6}$ W/m$^2$, intensitas gelombang gempa
itu pada jarak 400 km dari sumbernya adalah \ldots\ W/m$^2$.
\begin{pil}
  \item $2{,}50\times10^4$ \item $2{,}50\times10^5$ \item $6{,}25\times10^4$
  \item $6{,}25\times10^5$ \item $7{,}50\times10^5$
\end{pil}
\end{soal}

\begin{soal}{75}
Pada jarak 30 m dari suatu mesin jet yang beroperasi terukur intensitas
suara $1$ W/m$^2$. Intensitas suara yang tidak mengganggu percakapan
maksimum adalah $100\,\mu$W/m$^2$. Pada jarak berapa agar suara mesin
jet tidak mengganggu percakapan?
\begin{pil}
  \item minimum 3.000 m \item maksimum 3.000 m \item minimum 30.000 m
  \item maksimum 30.000 m \item minimum 300 m
\end{pil}
\end{soal}

% ===================================================================
%  SET 6  (nomor 76--90)
% ===================================================================

\begin{soal}{76}
Sebuah partikel bergerak pada sumbu $x$. Saat $t=0$ posisinya di
$x=+5$ m, saat $t=6$ s di $x=-7$ m, dan saat $t=10$ s di $x=+2$ m.
Kelajuan rata-rata partikel pada selang waktu antara $t=0$ sampai
$t=10$ s adalah \ldots\ m/s.
\begin{pilB}
  \item 0,3 \item 1,2 \item 2,1 \item 2,6 \item 2,7
\end{pilB}
\end{soal}

\begin{soal}{77}
Posisi $x$ sebuah partikel yang bergerak pada sumbu $x$ dinyatakan
dengan $x(t)=3t^2-12t+3$ (dalam meter). Pada selang waktu dari $t=0$ s
sampai $t=5$ s, kelajuan rata-rata partikel adalah \ldots\ m/s.
\begin{pilB}
  \item 2,4 \item 3,0 \item 4,8 \item 5,4 \item 7,8
\end{pilB}
\end{soal}

\begin{soal}{78}
Bola dipukul agar masuk ke parit $ab$ yang berada di lembah seperti pada
gambar. Sesaat setelah dipukul bola bergerak mendatar dengan laju $v_0$.
Kedalaman parit 3,2 m; tepi $a$ pada jarak mendatar 2,0 m dari posisi
bola dan lebar parit $ab=1{,}2$ m. Batas nilai $v_0$ agar bola masuk
ke parit $ab$ adalah \ldots
\begin{pil}
  \item $2{,}5<v_0<4{,}5$ m/s \item $2{,}5<v_0<4{,}0$ m/s
  \item $1{,}5<v_0<4{,}0$ m/s \item $3{,}2<v_0<4{,}2$ m/s
  \item $2{,}0<v_0<4{,}0$ m/s
\end{pil}
\end{soal}

\begin{soal}{79}
Sebuah bola kecil diikat dengan tali kemudian diputar dengan kecepatan
sudut yang tetap pada radius 0,25 m. Lintasan bola pada bidang horizontal
yang tingginya 1,25 m di atas lantai. Diketahui $g=10$ m/s$^2$. Pada
suatu saat tali putus, akibatnya bola jatuh pada jarak mendatar 2,0 m
dari posisi saat tali putus. Laju sudut bola saat berputar adalah
\ldots\ rad/s.
\begin{pilB}
  \item 16 \item 18 \item 20 \item 24 \item 25
\end{pilB}
\end{soal}

\begin{soal}{80}
Balon udara turun vertikal dengan laju tetap 20 m/s. Pada ketinggian
tertentu, dari balon dilemparkan sebuah batu secara mendatar dengan laju
15 m/s. Diamati bahwa batu menumbuk tanah enam detik sejak dilemparkan.
Perhatikan pernyataan: (1)~saat batu dilemparkan, posisi balon 300 m di
atas tanah; (2)~saat batu menyentuh tanah, posisi balon 120 m di atas
tanah; (3)~batu menumbuk tanah pada jarak mendatar 90 m dari posisi saat
dilemparkan. Pernyataan yang benar adalah \ldots
\begin{pilB}
  \item hanya (1) \item hanya (2) \item hanya (3)
  \item hanya (1) dan (3) \item (1), (2), dan (3)
\end{pilB}
\end{soal}

\begin{soal}{81}
Balon udara sedang bergerak naik vertikal dengan laju 10 m/s. Pada saat
balon 40 m di atas tanah, sebuah bola dilemparkan dari balon dengan
kecepatan awal $v_0$ dan arahnya tegak lurus terhadap lintasan balon.
Jika bola jatuh pada jarak mendatar 60 m dari garis lintasan balon,
maka $v_0=\ldots$ m/s.
\begin{pilB}
  \item 10 \item 15 \item 20 \item 22 \item 25
\end{pilB}
\end{soal}

\begin{soal}{82}
Peluru bermassa $m$ ditembakkan ke arah balok yang rehat dan bermassa
$M$ dan berada di tepi meja yang tingginya $h$. Peluru bersarang di
dalam balok, dan menyebabkan balok jatuh ke lantai pada jarak mendatar
$d$. Laju peluru saat menumbuk balok adalah \ldots
\gambar[0.45\textwidth]{images/2/2-p61.png}
\begin{pil}
  \item $\left(\dfrac{m+M}{M}\right)d\sqrt{\dfrac{g}{2h}}$
  \item $\left(\dfrac{M}{m}\right)d\sqrt{\dfrac{g}{2h}}$
  \item $\left(\dfrac{m+M}{m}\right)d\sqrt{\dfrac{g}{2h}}$
  \item $\left(\dfrac{m+M}{m}\right)d\sqrt{\dfrac{2g}{h}}$
  \item $\left(\dfrac{m}{M}\right)d\sqrt{\dfrac{g}{2h}}$
\end{pil}
\end{soal}

\begin{soal}{83}
Proyektil bermassa $m=10$ g ditembakkan ke atas dengan laju $v=1000$ m/s.
Proyektil itu menumbuk balok bermassa $M=5$ kg yang mula-mula rehat.
Proyektil menembus balok melalui pusat massa dan keluar dengan kecepatan
400 m/s menyebabkan balok naik setinggi maksimum $h$ dari posisi
awalnya. Jika $g=10$ m/s$^2$, maka $h=\ldots$
\begin{pilB}
  \item 5,0 cm \item 6,5 cm \item 7,0 cm \item 7,2 cm \item 7,5 cm
\end{pilB}
\end{soal}

\begin{soal}{84}
Balok bermassa $m_1$ dilepaskan di titik $A$ yang tingginya $h=2{,}7$ m.
Balok meluncur pada bidang lengkung yang licin kemudian menumbuk lenting
sempurna balok bermassa $m_2=2m_1$ yang awalnya diam di $B$. Akibat
tumbukan itu balok B meluncur pada bidang datar kasar yang koefisien
gesek kinetiknya 0,50. Balok $m_2$ meluncur pada bidang kasar sejauh
\ldots\ meter.
\gambar[0.45\textwidth]{images/2/2-p63.png}
\begin{pilB}
  \item 2,1 \item 2,2 \item 2,3 \item 2,4 \item 2,5
\end{pilB}
\end{soal}

\begin{soal}{85}
Gambar di bawah menunjukkan sebuah meja bundar yang berputar pada sumbu
vertikal dengan laju sudut tetap $\omega=2{,}5$ rad/s. Massa meja 25 kg
dan radiusnya $R=60$ cm. Sebuah balok kecil bermassa 200 g diletakkan
di meja pada jarak $r$ dari sumbu rotasi sedemikian sehingga balok
tidak mengalami slip terhadap meja. Jika koefisien gesek balok ke meja
0,25, maka $r=\ldots$ cm.
\gambar[0.5\textwidth]{images/2/2-p66.png}
\begin{pilB}
  \item 30 \item 35 \item 40 \item 45 \item 50
\end{pilB}
\end{soal}

\begin{soal}{86}
Sebuah mobil melintasi sebuah tikungan datar berupa busur lingkaran
berjari-jari 30 m. Besar koefisien gesek statik antara ban dan jalan
adalah $\mu_s=0{,}75$ dan $g=10$ m/s$^2$. Laju maksimum mobil melewati
tikungan agar tidak mengalami slip adalah \ldots\ m/s.
\begin{pilB}
  \item 15 \item 20 \item 25 \item 30 \item 40
\end{pilB}
\end{soal}

\begin{soal}{87}
Sebuah balok bermassa $m_1=2$ kg diikatkan pada tali yang panjangnya
$L_1=50$ cm dengan salah satu ujung tali terikat. Massa itu bergerak
dalam lintasan lingkaran horizontal pada meja yang licin. Balok kedua
bermassa $m_2=3$ kg diikatkan pada balok pertama dengan tali yang
panjangnya $L_2=50$ cm. Kedua balok berotasi bersama-sama dengan laju
sudut $\omega=5$ rad/s. Tegangan pada tali yang menghubungkan $m_1$ ke
pusat rotasi adalah \ldots\ N.
\begin{pilB}
  \item 25,0 \item 62,5 \item 75,0 \item 82,5 \item 100,0
\end{pilB}
\end{soal}

\begin{soal}{88}
Partikel bermassa $m$ terikat pada tali dan berputar horizontal di atas
meja yang kasar dengan kelajuan linier mula-mula $v_0$. Saat
menyelesaikan satu putaran penuh, kelajuan partikel turun menjadi
$\tfrac{2}{3}v_0$ akibat gaya gesek kinetik. Berapa putaran lagi yang
dapat dilakukan partikel sampai berhenti?
\begin{pilB}
  \item 0,6 \item 0,8 \item 1 \item 1,5 \item 2
\end{pilB}
\end{soal}

\begin{soal}{89}
Sebuah derek listrik harus menghasilkan energi 120 Wh saat digunakan
mengangkat balok besi 7200 N dari permukaan tanah tegak lurus ke atas.
Jika efisiensi kerja derek itu hanya 60\%, tinggi maksimum yang dicapai
balok besi dari permukaan tanah adalah \ldots\ m.
\begin{pilB}
  \item 24 \item 28 \item 30 \item 36 \item 48
\end{pilB}
\end{soal}

\begin{soal}{90}
Air terjun yang tingginya 100 m mengalir dengan debit 200 m$^3$/s. Di
dasar air terjun dipasang generator pembangkit listrik yang efisiensi
kerjanya 75\%. Besar percepatan gravitasi setempat 10 m/s$^2$. Daya
keluaran generator tersebut adalah \ldots\ MW.
\begin{pilB}
  \item 100 \item 150 \item 200 \item 250 \item 300
\end{pilB}
\end{soal}

\end{document}
